\documentclass[11pt]{article}
\usepackage[T1]{fontenc}
\usepackage[utf8]{inputenc}
\usepackage{lmodern}
\usepackage{amsmath,amssymb,bm}
\usepackage{geometry}
\usepackage{hyperref}
\geometry{margin=1in}

\title{Pion Vector-Current Matrix Elements in \texttt{jaxQFT}}
\author{}
\date{\today}

\begin{document}
\maketitle

\section{Scope}
This note documents the DD-independent analysis path used to extract a pion
vector-current matrix element from two-point and three-point functions in the
two-dimensional U(1) Wilson-fermion model used in this repository as a
Schwinger-model proxy.

The goal is narrow and practical:
\begin{enumerate}
  \item fit the pion two-point functions to obtain the single-particle energy
  $E_\pi(p)$ and the overlap factor $Z_\pi(p)$,
  \item use those fits together with a pion three-point function to remove the
  source/sink overlap factors,
  \item extract the matrix element
  $\langle \pi(p_f)|V_\mu|\pi(p_i)\rangle$,
  \item for the temporal current, convert that matrix element to the pion form
  factor.
\end{enumerate}

\section{Current Choice}
For an isospin-symmetric two-flavor theory the correct connected current for the
charged-pion form factor is the isovector vector current
\begin{equation}
  V_\mu^{(3)}(x)
  =
  \bar\psi(x)\,\gamma_\mu\,\frac{\tau^3}{2}\,\psi(x).
\end{equation}
If the exactly conserved point-split Wilson current is used, its renormalization
constant is
\begin{equation}
  Z_V = 1.
\end{equation}
If instead a local current is used, the analysis stays the same but a separate
current renormalization factor must be applied at the end. The script therefore
keeps a user-visible multiplicative input $Z_V$ and defaults it to $1$.

\section{Single-Pion Two-Point Function}
For a local pseudoscalar interpolator
\begin{equation}
  P^a(x) = \bar\psi(x)\,\gamma_5\,\tau^a\,\psi(x),
\end{equation}
the momentum-projected two-point function is
\begin{equation}
  C_2(t;\,p)
  =
  \sum_x e^{ipx}\,
  \langle P^a(x,t)\,P^a(0)^\dagger \rangle.
\end{equation}
At large Euclidean time,
\begin{equation}
  C_2(t;\,p)
  \approx
  \frac{|Z_\pi(p)|^2}{2E_\pi(p)}
  \left(
    e^{-E_\pi(p)t}
    +
    e^{-E_\pi(p)(T-t)}
  \right),
\end{equation}
with
\begin{equation}
  Z_\pi(p) = \langle 0|P^a(0)|\pi^a(p)\rangle.
\end{equation}

In the code this is fit as
\begin{equation}
  C_2(t;\,p) = A(p)\left(e^{-Et}+e^{-E(T-t)}\right),
\end{equation}
so the overlap factor used by the three-point analysis is
\begin{equation}
  Z_\pi(p) = \sqrt{2E_\pi(p)\,A(p)}.
\end{equation}

\section{Three-Point Function}
For source at time $0$, current insertion at time $\tau$, and sink at
$t_{\rm sep}$, define
\begin{equation}
  C_{3,\mu}(\tau,t_{\rm sep};p_f,p_i)
  =
  \sum_{x,y}
  e^{-ip_f x}
  e^{+i(p_f-p_i)y}
  \langle
    P(x,t_{\rm sep})\,
    V_\mu(y,\tau)\,
    P(0)^\dagger
  \rangle.
\end{equation}
In the single-state region $0 \ll \tau \ll t_{\rm sep}$,
\begin{equation}
  C_{3,\mu}(\tau,t_{\rm sep};p_f,p_i)
  \approx
  \frac{Z_\pi(p_f)\,Z_\pi(p_i)}{4E_fE_i}
  \langle \pi(p_f)|V_\mu|\pi(p_i)\rangle\,
  e^{-E_f(t_{\rm sep}-\tau)}\,e^{-E_i\tau},
\end{equation}
where $E_i \equiv E_\pi(p_i)$ and $E_f \equiv E_\pi(p_f)$.

\section{Peeling Off the Overlap Factors}
\subsection{Direct overlap removal}
Once $E_i$, $E_f$, $Z_\pi(p_i)$, and $Z_\pi(p_f)$ are known from the two-point
fits, the matrix element can be isolated directly:
\begin{equation}
  M_\mu^{(Z)}(\tau)
  \equiv
  Z_V\,
  \frac{4E_fE_i}{Z_\pi(p_f)\,Z_\pi(p_i)}
  e^{E_f(t_{\rm sep}-\tau)+E_i\tau}
  C_{3,\mu}(\tau,t_{\rm sep};p_f,p_i).
\end{equation}
In the plateau region,
\begin{equation}
  M_\mu^{(Z)}(\tau)
  \to
  Z_V\,\langle \pi(p_f)|V_\mu|\pi(p_i)\rangle.
\end{equation}

\subsection{Ratio method}
The same matrix element can be extracted with the standard ratio
\begin{equation}
  R_\mu(\tau,t_{\rm sep};p_f,p_i)
  =
  \frac{C_{3,\mu}(\tau,t_{\rm sep};p_f,p_i)}{C_2(t_{\rm sep};p_f)}
  \sqrt{
    \frac{
      C_2(t_{\rm sep}-\tau;p_i)\,
      C_2(\tau;p_f)\,
      C_2(t_{\rm sep};p_f)
    }{
      C_2(t_{\rm sep}-\tau;p_f)\,
      C_2(\tau;p_i)\,
      C_2(t_{\rm sep};p_i)
    }
  }.
\end{equation}
Using the same single-state asymptotics as above, one finds
\begin{equation}
  R_\mu(\tau,t_{\rm sep};p_f,p_i)
  \to
  \frac{
    \langle \pi(p_f)|V_\mu|\pi(p_i)\rangle
  }{
    2\sqrt{E_fE_i}
  }.
\end{equation}
The corresponding effective matrix element is therefore
\begin{equation}
  M_\mu^{(R)}(\tau)
  \equiv
  Z_V\,2\sqrt{E_fE_i}\,
  R_\mu(\tau,t_{\rm sep};p_f,p_i).
\end{equation}

This ratio is useful because the overlap factors cancel explicitly. In the
current implementation the plateau window is selected from $M_\mu^{(R)}(\tau)$,
and the same window is then applied to both $M_\mu^{(R)}$ and $M_\mu^{(Z)}$ as
a cross-check.

On the small lattices used for this Schwinger-model workflow, the two-point
function can still have visible backward propagation even when the three-point
function is analyzed in a forward-state window. For that reason, the code uses
the fitted forward branch
\begin{equation}
  C_2^{\rm fw}(t;\,p) = A(p)e^{-E_\pi(p)t}
\end{equation}
inside the ratio estimator by default. This keeps the ratio and direct
overlap-removal estimators on the same single-state footing.

\section{Form-Factor Normalization}
For the temporal current in Euclidean kinematics,
\begin{equation}
  \langle \pi(p_f)|V_t|\pi(p_i)\rangle
  =
  \left(E_f + E_i\right) F_\pi(Q^2),
\end{equation}
with
\begin{equation}
  Q^2 = (E_f-E_i)^2 + (p_f-p_i)^2
\end{equation}
in the continuum-like convention used by the script. Therefore
\begin{equation}
  F_\pi(Q^2)
  =
  \frac{
    \langle \pi(p_f)|V_t|\pi(p_i)\rangle
  }{
    E_f + E_i
  }.
\end{equation}
The script reports this quantity automatically when the selected current index
is the temporal one.

\section{Statistical Procedure}
The extraction code is designed to stay consistent with the rest of the MCMC
analysis pipeline:
\begin{enumerate}
  \item Align the two-point and three-point records on the same MCMC steps.
  \item Estimate an integrated autocorrelation time from the two-point and
  three-point data and choose a block size $b \gtrsim 2\tau_{\rm int}$ unless
  the user overrides it.
  \item Build blocked means.
  \item Fit the pion two-point functions on blocked data and propagate the
  resulting $E_\pi(p)$ and $Z_\pi(p)$ through a blocked jackknife.
  \item Construct the effective matrix-element curves on each jackknife sample.
  \item Choose the plateau window from a correlated constant fit to the ratio
  estimator.
  \item Quote the final value and error from the same jackknife samples and the
  selected window.
\end{enumerate}

\section{Expected Three-Point Record Layout}
The current analysis script reads checkpoint records with flat scalar keys of
the form
\begin{verbatim}
c3_mu{mu}_pf{pf}_pi{pi}_tsep{tsep}_tau{tau}_re
c3_mu{mu}_pf{pf}_pi{pi}_tsep{tsep}_tau{tau}_im
\end{verbatim}
inside a measurement block whose name defaults to
\texttt{pion\_3pt\_vector}. This fixes the interface for the future inline
three-point measurement implementation.

\section{Code Path}
The extraction workflow is implemented in
\begin{verbatim}
scripts/mcmc/analyze_pion_vector_3pt.py
\end{verbatim}
and uses the existing correlated two-point fitting machinery from
\texttt{scripts/mcmc/fit\_2pt\_dispersion.py}. The result is a JSON summary and,
optionally, a PNG showing the effective matrix-element plateau.

\section{Practical Caveat}
The direct overlap-removal estimator assumes the forward single-state form of
the three-point function. In practice this is safe only when
\begin{equation}
  0 \ll \tau \ll t_{\rm sep} \ll T/2.
\end{equation}
If wrap-around effects or excited states are significant, the plateau will not
be reliable and one must move to a more complete simultaneous
two-state/two-point/three-point fit.

\end{document}
